\documentclass[../main.tex]{subfiles}

\begin{document}

\chapter{$L^{p}$ 空间: 完备性与逼近}\label{chapter:ch07-Lp空间}

实数的完备性可阐述为: 若 $\{a_{n}\}$ 是满足 $\lim_{n,m\to\infty}|a_{n}-a_{m}|=0$ 的实数序列, 则存在实数 $a$ 使得 $\lim_{n\to\infty}|a_{n}-a|=0$. 对于 Lebesgue 积分有与之对应的完备性. 对于可测集 $E$ 和 $1\leqslant p<\infty$, 定义 $L^{p}(E)$ 为满足 $|f|^{p}$ 在 $E$ 上可积的可测函数族, 因此 $L^{1}(E)$ 是可积函数族. 若 $\{f_{n}\}$ 是 $L^{p}(E)$ 中满足
\begin{align*}
\lim_{n,m\to\infty}\int_{E}|f_{n}-f_{m}|^{p}=0
\end{align*}
的函数序列, 则存在属于 $L^{p}(E)$ 的函数 $f$, 使得
\begin{align*}
\lim_{n\to\infty}\int_{E}|f_{n}-f|^{p}=0.
\end{align*}
这是 Riesz-Fischer 定理——本章的中心. $L^{p}(E)$ 中的函数族 $\mathcal{F}$ 称为在 $L^{p}(E)$ 中稠密, 若对每个 $L^{p}(E)$ 中的函数 $g$ 和 $\varepsilon>0$, 存在属于 $\mathcal{F}$ 的函数 $f$, 使得 $\int_{E}|g-f|^{p}<\varepsilon$. 我们证明在 $L^{p}(E)$ 中存在稠密的可数函数族, 且连续函数与简单函数都在 $L^{p}(E)$ 中稠密. Riesz-Fischer 定理与稠密性结果的证明在函数的赋范线性空间的背景下构建. 为构造这个框架, 我们证明两个基本的不等式——H\"older 不等式与 Minkowski 不等式.

\subfile{Chapter-07/section-01.tex}

\subfile{Chapter-07/section-02.tex}

\subfile{Chapter-07/section-03.tex}

\subfile{Chapter-07/section-04.tex}

\end{document}